\section{Introduction}

Self-evolving agents adapt without updating model weights by converting execution experience into persistent procedural state. Recent systems already show that failed trajectories, success/failure contrasts, execution-grounded repair, and accumulated knowledge can improve later skill use~\citep{liu2026rethinkskill,chen2026skillcat,liu2026skillrevise,tang2026wikiskill}. Most evaluations, however, emphasize \emph{which evidence} reaches the updater. They largely treat the state-writing step itself as an implementation detail.

We study a narrower failure mode: the same learner-visible evidence can fail to regenerate the same behaviorally useful persistent skill. This matters because persistent skill text is itself an intervention on future behavior. Two updater calls can receive identical evidence yet synthesize states that differ in scope, redundancy, ordering, specificity, or procedural commitments. Thus evidence selection and state materialization are experimentally separable even when they are implemented by the same LLM pipeline.

Our starting point is a controlled DeepSeek-V4-Pro study of persistent spreadsheet skills. A 48-pair experiment shows that replacing winner evidence with rejected-witness evidence has only a small heterogeneous average effect, so the paper cannot be about universally superior failure feedback. In one outcome-selected development stream, however, a historically strong First-Fail state remains directionally useful when its bytes are frozen and re-evaluated. Reconstructing the learner-visible First-Fail evidence byte-for-byte and invoking the same free-form updater twice does \emph{not} recreate that historical advantage. The evidence remained fixed; the useful state did not reliably reappear.

We organize the pipeline as
\begin{equation}
\text{trajectories}\rightarrow E\rightarrow G\rightarrow S\rightarrow Y,
\end{equation}
where $E$ selects the learner-visible evidence package, $G$ generates persistent state, $S$ is the realized skill artifact, and $Y$ is future utility. This factorization is bookkeeping rather than a standalone novelty claim. Its purpose is to make the realized persistent state an explicit experimental object. The completed evidence establishes only a \emph{local state-regeneration instability}: one selected case in which identical evidence does not reliably regenerate the historical useful state. It does not estimate a population variance component or show that updater variance dominates actor noise.

This observation motivates a controlled intervention on $G$. We replace unconstrained state writing with a small typed compiler that maps the same trajectory+score package to a finite diagnosis/repair representation and a canonical skill state. The compiler is not presented as a universal diagnosis algorithm: SkillRevise, SkillCAT, WikiSkill, and related systems already cover substantial parts of trajectory diagnosis, contrastive feedback, and persistent knowledge construction~\citep{liu2026skillrevise,chen2026skillcat,tang2026wikiskill}. Its role here is causal control: remove free-form state-synthesis degrees of freedom while preserving learner-visible evidence.

The prospective bridge therefore separates four questions that earlier versions conflated. The primary question is the \emph{generator factor}: averaged equally across predeclared Winner and First-Fail-4 evidence cells, does constrained generation outperform the native free-form updater? A second, explicitly FF4-specific question asks whether the compiled state survives score-only and diagnosis-cardinality-informed trajectory-text-blind controls; this classifies the substance of the FF4 typed-diagnosis result and cannot explain the Winner-side contribution to the balanced main effect. A third same-evidence repeated-synthesis probe asks whether distinct FF4 FREE state realizations show observed cross-state disagreement beyond re-running the actor on frozen states; a stronger success-propensity interpretation is allowed only under an explicit conditionally iid/independent stationary actor model. First-Fail superiority and the Evidence$\times$Generator interaction are a fourth, parked moderator question. These gates are orthogonal: failure of a secondary interpretation cannot manufacture or erase the primary generator-factor result.

\paragraph{Contributions.}
\begin{itemize}
    \item \textbf{Matched-evidence regeneration audit.} In one controlled selected case, a historically useful state remains behaviorally active when frozen, while byte-identical evidence does not reliably regenerate its advantage through the native updater.
    \item \textbf{Bounded causal localization.} Together with a 48-pair heterogeneous evidence study and a failed witness-selector pilot, this rules out evidence disappearance and weakens a one-rollout actor-luck explanation while explicitly leaving selected-case bias and population variance unresolved.
    \item \textbf{Controlled generator intervention and falsification design.} We define a typed canonical compiler, strong trajectory-blind generic controls, SHA-based state-identity handling, and a fresh balanced Evidence$\times$Generator bridge whose primary generator-factor gate is separated from diagnosis, realization, and rejected-source interpretation. Its STOP rules prevent E3, another backbone, or another benchmark from rescuing a failed generator result.
\end{itemize}
